The future works of our {path-sensitive} reachability-bound analysis
mainly focus on improving our algorithm and the implementation
from the three aspects,
% has mainly three limitations in terms of
the \emph{efficiency}, \emph{accuracy} and the \emph{generality}.

\subsubsection{The Efficiency} 
% The efficiency of the algorithm.
In our algorithm, there are six steps in order to precisely estimate the reaching times of each location. Each step involves at least one iteration of the program, which results in low efficiency.
The major bottleneck comes from the recursion and potential non-termination in the path-refinement algorithm and computing the invariants for the ranking functions. So in the next step, we will explore new path refinement~\cite{CyphertBKR19} and loop summarization techniques~\cite{BreckCKR20,KincaidBCR19,KincaidCBR18} and design a more efficient path refinement algorithm. We also aim to design a more efficient ranking estimation in replace of the existing one.

\subsubsection{The Accuracy} 
The accuracy of the path refinement algorithm and program abstraction implementation.
We sacrificed accuracy in order to improve the efficiency of the implementation. For the path refinement algorithm,
we limited its maximal iteration times to $100$, and the same for computing program constraints. 
In the future, we plan to improve the accuracy of the implementation by extending the implementation with the new efficient algorithm.
\cite{CyphertBKR19} introduced a precise path refinement technique over regular expression. Its experiments show that they are accurate in eliminating infeasible loop paths, which can be a potential technique for improvement.

\subsubsection{The Generality}
In our existing implementation, the prototype is limited in terms of expressiveness and scalability. So far we only support the program with integer counters and linear counter transformation. While there are some advanced works reasoning about the non-linear loop summarization and invariant generation. Against this background, we plan to support more complex computations in the next step.

